% Drafted Introduction + Related Work for Stage 3 to ingest.
% All citations resolve to entries in references.bib.

\section{Introduction}

Continuous learning and memory formation across the lifespan depend on both developmental and adult neurogenesis in the dentate gyrus (DG) of the hippocampus, one of the principal neurogenic niches in the mammalian brain. In rodents, DG neural stem cells (NSCs) transition reversibly between quiescent and active states, sustain granule-cell production into adulthood, and can be reactivated by environmental or pathological cues \cite{encinas2011division,obernier2019neural,urban2019quiescence}. Whether a comparable neurogenic process persists in the adult human hippocampus has been debated for more than two decades. Traditional immunostaining studies have produced diametrically opposed conclusions: some report sharply declining numbers of doublecortin-positive (DCX$^{+}$) cells that become undetectable in adults \cite{sorrells2018human,cipriani2018hippocampal,franjic2022transcriptomic}, whereas others document persistent adult neurogenesis and an influence of disease state or exercise on the neurogenic pool \cite{boldrini2018human,morenojimenez2019adult,tobin2019human,terrerosroncal2021impact}. Reconciling these claims is difficult because human tissue is scarce, postmortem quality is variable, and the molecular markers most commonly used in rodents are not necessarily specific in the human hippocampus.

Single-nucleus RNA sequencing (snRNA-seq) offers a path around several of these limitations: it operates on frozen postmortem material, avoids antibody-specificity biases, and resolves rare cell populations that are invisible in bulk assays \cite{habib2017massively}. Large-scale snRNA-seq of the developing \cite{zhong2020decoding} and adult \cite{franjic2022transcriptomic,ayhan2021resolving} human hippocampus has begun to catalogue its cell types, yet the molecular dynamics of NSCs themselves --- how quiescent, primed, and active states evolve across life stages, and whether injury can reawaken them in humans --- remain unresolved. In rodents, dormant subventricular-zone NSCs re-enter the cell cycle after ischemia through IFN$\gamma$-primed states \cite{llorensbobadilla2015single}; whether equivalent injury-driven reactivation operates in the human hippocampus has not been tested with transcriptomic resolution.

Here we generate a single-nucleus transcriptomic atlas of the anterior-mid human hippocampus spanning neonatal (postnatal day 4), adult (31--32 years), aging (50--68 years), and stroke-injured (48 years) states, comprising 92{,}966 quality-controlled nuclei from 10 donors. Using cross-species integration with mouse, pig, and macaque data \cite{hochgerner2018conserved,franjic2022transcriptomic}, single-cell hierarchical Poisson factorization \cite{levitin2019gene}, Seurat-based marker discovery and integration \cite{butler2018integrating}, RNA-velocity \cite{lamanno2018rna}, and Monocle~2 pseudotime \cite{trapnell2014dynamics}, we (i) resolve the NSC lineage into quiescent (qNSC1, qNSC2), primed (pNSC), active (aNSC), and neuroblast (NB) populations; (ii) identify human-specific markers that avoid the GABAergic-interneuron contamination confounding DCX, PROX1, and STMN2; (iii) show that most NSCs enter a deep-quiescence state during aging; and (iv) demonstrate that stroke injury reactivates quiescent NSCs into a proliferative primed/active state that recapitulates neonatal neurogenic signatures. Together, these data offer a concrete molecular basis for why adult human hippocampal neurogenesis is rarely observed under healthy conditions but can be detected in brains under specific pathological pressure \cite{morenojimenez2019adult,terrerosroncal2021impact}.

\paragraph{Contributions.}
\begin{itemize}
    \item A single-nucleus transcriptomic atlas of the human hippocampus covering neonatal, adult, aging, and stroke-injured donors, with 92{,}966 quality-controlled nuclei partitioned into 16 major populations.
    \item Cross-species-validated subdivision of the AS2/qNSC compartment into astrocyte2, qNSC1 (deep quiescence, HOPX-high), and qNSC2 (shallow quiescence, LPAR1/STMN1-high), with molecular signatures and cell-cycle properties distinguishing each from astrocytes and each other.
    \item Human-specific neurogenic-lineage markers --- LRRC3B, RHOJ, SLC4A4, GLI3 (qNSC); CHI3L1, EGFR (pNSC); and CALM3, TTC9B, NRGN, FXYD7, NRN1, GNG3, TCEAL5, TMSB10, NEUROD2 (NB) --- identified by combining scHPF and Seurat FindAllMarkers and filtered against GABAergic-interneuron expression to avoid the interneuron contamination of canonical markers.
    \item A developmental trajectory of the neonatal DG NSC lineage supported by RNA-velocity and Monocle~2 pseudotime, with gene-expression and transcription-factor cascades from pNSCs to two granule-cell neuron subtypes, validated by HOPX, NES, Ki67, and PSA-NCAM immunostaining.
    \item Quantitative evidence that the neurogenic lineage shifts toward deep quiescence during aging and is partially reactivated by stroke, with pNSC and aNSC fractions recovering to 17.3\% (322/1861) and 11.7\% (218/1861) of the injured-hippocampus neurogenic lineage --- comparable to neonatal ratios and markedly higher than in adult and aging groups.
\end{itemize}

\section{Related Work}

\paragraph{The human adult hippocampal neurogenesis debate.}
Early immunohistochemical studies in adult human hippocampus reached opposing conclusions. Sorrells et al.\ reported that DCX$^{+}$ young neurons are essentially undetectable after the first year of life \cite{sorrells2018human}, and Cipriani et al.\ found that dentate radial-glia proliferation is negligible in the adult DG \cite{cipriani2018hippocampal}. In contrast, Boldrini et al.\ quantified thousands of immature neurons into the eighth decade \cite{boldrini2018human}, Moreno-Jim\'{e}nez et al.\ reported abundant DCX$^{+}$ cells in neurologically healthy donors that sharply decline in Alzheimer's disease \cite{morenojimenez2019adult}, and Tobin et al.\ detected persistent neurogenesis in cognitively impaired ninth-decade brains \cite{tobin2019human}. Terreros-Roncal et al.\ extended the latter view by showing neurodegeneration-associated alterations in the neurogenic niche across multiple diseases \cite{terrerosroncal2021impact}. Our data indicate that these apparently contradictory observations can be unified once the deep-quiescence state and its reactivation by injury or disease are taken into account.

\paragraph{Cross-species scRNA-seq of the hippocampal neurogenic lineage.}
Single-cell profiling of the mouse DG has established a continuous trajectory from RGLs through intermediate progenitors to mature granule cells with conserved marker sets \cite{hochgerner2018conserved}, and similar single-cell analyses of the mouse subventricular zone have mapped the quiescence--activation axis \cite{dulken2017single,llorensbobadilla2015single}. In primates, Franjic et al.\ combined snRNA-seq of human, macaque, and pig hippocampal--entorhinal tissue and reported robust neurogenic signatures in non-human species but not in adult humans, additionally demonstrating that DCX is broadly expressed in diverse human neurons rather than being restricted to newly generated granule cells \cite{franjic2022transcriptomic}. Ayhan et al.\ further characterized anterior/posterior molecular diversity in surgically resected adult human hippocampi \cite{ayhan2021resolving}, while Zhong et al.\ produced a developmental atlas from 16--27 gestational weeks anchored by PAX6$^{+}$ and HOPX$^{+}$ progenitors \cite{zhong2020decoding}. Our atlas complements these studies by spanning neonatal to aged donors and by including an injury condition absent from prior work.

\paragraph{NSC quiescence, activation, and injury response.}
Rodent work has long supported a quiescent NSC pool that can be mobilized. Encinas et al.\ showed that dividing hippocampal NSCs ultimately convert into astrocytes, depleting the pool with age \cite{encinas2011division}. Obernier and Alvarez-Buylla reviewed the heterogeneity and regulatory inputs that define adult NSCs \cite{obernier2019neural}, and Urban et al.\ synthesized the molecular logic of NSC quiescence \cite{urban2019quiescence}. At the single-cell level, Llorens-Bobadilla et al.\ identified dormant and primed-quiescent subpopulations that activate upon brain ischemia through IFN$\gamma$-linked signalling \cite{llorensbobadilla2015single}, and Dulken et al.\ resolved the adult NSC-to-neuroblast continuum by scRNA-seq \cite{dulken2017single}. Whether analogous injury-induced reactivation exists in the human hippocampus had not been tested; our stroke-donor sample provides direct transcriptomic evidence that it does.

\paragraph{Reactive astrocytes and NSC ambiguity.}
Because human qNSCs and astrocytes share GFAP, ALDH1L1, and VIM, distinguishing them transcriptomically requires careful gene-set scoring. Zamanian et al.\ catalogued injury-induced reactive astrocyte transcriptomes \cite{zamanian2012genomic}; Liddelow et al.\ defined neurotoxic A1 astrocytes and their inducers \cite{liddelow2017neurotoxic}; and Clarke et al.\ showed that normal aging itself drives A1-like reactivity \cite{clarke2018normal}. Our analysis uses these gene sets to separate reactive astrocytes from reactivated NSCs in the stroke-injured hippocampus, a distinction essential for interpreting the apparent rebound of pNSC- and aNSC-compatible transcriptomes after injury.

\paragraph{Computational methods for lineage and marker discovery.}
To resolve rare populations and distinguish transcriptionally similar cell types, we combine several established methods. Seurat's integration framework harmonizes datasets across conditions and species \cite{butler2018integrating}; scHPF performs sparsity-aware factorization that recovers specific rather than merely abundant gene signatures \cite{levitin2019gene}; Monocle~2 orders cells along inferred trajectories through DDRTree \cite{trapnell2014dynamics}; and RNA velocity leverages spliced/unspliced transcript ratios to estimate directional lineage flow \cite{lamanno2018rna}. DroNc-seq provides the technical substrate by enabling snRNA-seq from postmortem frozen tissue \cite{habib2017massively}. Our contribution is to apply these tools jointly to a carefully staged human cohort --- neonatal, adult, aging, and injured --- to recover both baseline and injury-evoked NSC dynamics that cannot be seen in any single prior dataset, and to cross-reference findings with primate-specific observations such as the protracted $>$6-month granule-cell maturation of macaques \cite{kohler2011maturation} that underlies the differential longevity of human neuroblasts.
